\documentclass[memoire, 12pt]{report}
\usepackage[top = 1.9cm, bottom = 1.5cm, left = 1.9cm, right = 2.1cm]{geometry}
\usepackage{graphicx} 
\usepackage{enumitem}
\usepackage{multicol}
\usepackage{tabto}
\usepackage{multirow}
\usepackage{multibib}
\usepackage{multirow}
\usepackage{tabularx}
\newcites{biblio}{Bibliographie}
\newcites{other}{Autres r\'ef\'erences}
\usepackage{amssymb}
\usepackage{amsmath}
\usepackage{graphicx}
\usepackage{amsfonts}
\usepackage{lmodern}
\usepackage{caption}
\usepackage{subcaption}
\usepackage[babel=true]{csquotes}
\setlength{\fboxrule}{0.01cm}
\setlength{\fboxsep}{0.5cm}
\usepackage{array}
\usepackage{tikz}
\usepackage{lipsum}
\usepackage{setspace}
\usepackage{ragged2e}
\usepackage{url}
\usepackage{float}
\usepackage{pdfpages}
\usepackage{rotating}
\usepackage{glossaries}
%\usepackage[thinlines]{easytable}
\usepackage{hyperref}
\usepackage[export]{adjustbox}
\usepackage[bottom]{footmisc}
%\usepackage{algpseudocode}
\usepackage{algorithm}
\usepackage{algorithmic}
\usepackage[normalem]{ulem}
\useunder{\uline}{\ul}{}
\usepackage{glossaries}
\usepackage{listings}
\usepackage{xcolor}
\usepackage{minted}

\usepackage{array}
\usepackage{longtable}
\usepackage[table,xcdraw]{xcolor}

\usepackage[utf8]{inputenc}   % encodage du fichier source
\usepackage[T1]{fontenc}      % encodage des polices
\usepackage[french]{babel}    % pour le français

% Configuration des styles pour le code Python

\definecolor{codegreen}{rgb}{0,0.6,0}
\definecolor{codegray}{rgb}{0.5,0.5,0.5}
\definecolor{codepurple}{rgb}{0.58,0,0.82}
\definecolor{backcolour}{rgb}{0.95,0.95,0.92}

\lstdefinestyle{python}{
    backgroundcolor=\color{backcolour},   
    commentstyle=\color{codegreen},
    keywordstyle=\color{magenta},
    numberstyle=\tiny\color{codegray},
    stringstyle=\color{codepurple},
    basicstyle=\ttfamily\footnotesize,
    breakatwhitespace=false,         
    breaklines=true,                 
    captionpos=b,                    
    keepspaces=true,                 
    numbers=left,                    
    numbersep=5pt,                  
    showspaces=false,                
    showstringspaces=false,
    showtabs=false,                  
    tabsize=2
}

\lstset{style=python}



\usepackage[utf8]{inputenc}  
\usepackage[T1]{fontenc} 
%\usepackage{fancyhdr}
\usepackage[Conny]{fncychap}
%Conny
%Bjornstrup
%\pagestyle{Conny}
\usepackage[french]{babel}
%\renewcommand{\footrulewidth}{3pt}
\makeglossaries
\title{TAF_N1_NGWAMBE MARIELLA}
\author{}
\date{MOIS_ICI 2025}

\begin{document}
\begin{titlepage}

	\begin{tikzpicture}[remember picture,overlay,inner sep=0,outer sep=0]
		\draw[orange!90!orange,line width=4pt] ([xshift=-1.5cm,yshift=-2cm]current page.north east) coordinate (A)--([xshift=1.5cm,yshift=-2cm]current page.north west) coordinate(B)--([xshift=1.5cm,yshift=2cm]current page.south west) coordinate (C)--([xshift=-1.5cm,yshift=2cm]current page.south east) coordinate(D)--cycle;
		
		\draw ([yshift=0.5cm,xshift=-0.5cm]A)-- ([yshift=0.5cm,xshift=0.5cm]B)--
		([yshift=-0.5cm,xshift=0.5cm]B) --([yshift=-0.5cm,xshift=-0.5cm]B)--([yshift=0.5cm,xshift=-0.5cm]C)--([yshift=0.5cm,xshift=0.5cm]C)--([yshift=-0.5cm,xshift=0.5cm]C)-- ([yshift=-0.5cm,xshift=-0.5cm]D)--([yshift=0.5cm,xshift=-0.5cm]D)--([yshift=0.5cm,xshift=0.5cm]D)--([yshift=-0.5cm,xshift=0.5cm]A)--([yshift=-0.5cm,xshift=-0.5cm]A)--([yshift=0.5cm,xshift=-0.5cm]A);
		
		
		\draw ([yshift=-0.3cm,xshift=0.3cm]A)-- ([yshift=-0.3cm,xshift=-0.3cm]B)--
		([yshift=0.3cm,xshift=-0.3cm]B) --([yshift=0.3cm,xshift=0.3cm]B)--([yshift=-0.3cm,xshift=0.3cm]C)--([yshift=-0.3cm,xshift=-0.3cm]C)--([yshift=0.3cm,xshift=-0.3cm]C)-- ([yshift=0.3cm,xshift=0.3cm]D)--([yshift=-0.3cm,xshift=0.3cm]D)--([yshift=-0.3cm,xshift=-0.3cm]D)--([yshift=0.3cm,xshift=-0.3cm]A)--([yshift=0.3cm,xshift=0.3cm]A)--([yshift=-0.3cm,xshift=0.3cm]A);

	\end{tikzpicture}
	\begin{center}
		\begin{tabular}{l*{40}{@{\hskip.05mm}c@{\hskip.8mm}} c c}
			\begin{tabular}{c}
				
		\footnotesize{\textbf{R\'EPUBLIQUE DU CAMEROUN}} \\
				
				\scriptsize{\textbf{****************}} \\
				
					\scriptsize{\textbf{Paix - Travail - Patrie}} \\
				
			\scriptsize{\textbf{******************}}\\ 
			\footnotesize{	\textbf{UNIVERSIT\'E DE YAOUND\'E I}}\\
				
			\scriptsize{	\textbf{****************}} \\
				
			\footnotesize{	\textbf{ECOLE NATIONALE SUPERIEURE}} \\
			\footnotesize{	\textbf{POLYTECHNIQUE DE YAOUNDE}} \\
				
			\scriptsize{	\textbf{****************}} \\
		   \scriptsize{	\textbf{D\'EPARTEMENT DE GENIE}}\\
		   \scriptsize{	\textbf{INFORMATIQUE}}\\
				
			\scriptsize{	\textbf{****************}}\\
				
			\end{tabular} &
			\begin{tabular}{c}
				
				\includegraphics[height=4cm, width=2.8cm]{po.png}
				
			\end{tabular} &
			\begin{tabular}{c}
				
				\footnotesize{\textbf{ REPUBLIC OF CAMEROON}} \\
				
				\footnotesize{\textbf{****************}} \\
				
					\scriptsize{\textbf{Peace - Work - Fatherland}} \\
				
				\scriptsize{\textbf{****************}} \\
				\footnotesize{\textbf{UNIVERSITY OF YAOUNDE I}}\\
				
				\scriptsize{\textbf{****************}} \\
				
				\footnotesize{\textbf{NATIONAL ADVANCED SCHOOL}} \\
				\footnotesize{\textbf{OF ENGINEERING OF YAOUNDE}} \\
				
				\scriptsize{\textbf{****************}} \\
				\scriptsize{\textbf{DEPARTMENT OF COMPUTER}}\\
				\scriptsize{\textbf{ENGINEERING}}\\
				
				\footnotesize{\textbf{****************}}\\
				
			\end{tabular}	
		\end{tabular}
	
		\vspace{0.5cm}
		\begin{tabular}{l*{40}{@{\hskip 3.5cm}c@{\hskip5cm}} p{3.5cm} r}
		\end{tabular}
		
		\noindent\rule{\textwidth}{0.7mm}
		\Large{{\textbf{CHAPITRE 1}}}\\
		\Large{{\textbf{\textit{EXERCICES}}}}
		\noindent\rule{\textwidth}{0.7mm}
	\end{center}
		
	\begin{center}
	\begin{tabular}{c}
		
		\vspace{0.1cm}
		\normalsize
	
	
		\vspace{0.1cm}
		\normalsize\textbf{Option }:\\			
		\textsl{Cybersécurité et Investigation Numérique}
		
	\end{tabular}
	\end{center}
		
	\begin{center}
		\normalsize %\hspace{-2cm}
		\begin{tabular}{c}
			\vspace{0.07cm}
			\hspace{0.02cm} \textbf{\textbf{Rédigé par :}}\\
			\hspace{0.02cm} \textsl{\textbf{NGWAMBE  MARIELLA}, 22P040}\\\\
			
		\end{tabular}
	\end{center}
	
	\begin{center}
	\hspace{0.02cm} \textbf{Sous l'encadrement de:}\\
	\hspace{0.02cm} \textsl{M. Thierry MINKA}
	\end{center}
	
    
	\vspace{2cm}
	\begin{center}
		\textbf{Année académique 2025 / 2026}
	\end{center}
		
	\vspace{-1.4cm}
	
		
	\vfill%\null
	
\end{titlepage}
\begin{document}
 \chapter{} PARTIE 1: 
   \section{ Exercice 1}
 \paragraph{} — Paradoxe de la transparence}
La « société de la transparence » décrite par Byung-Chul Han met en lumière
une tension essentielle : la volonté de tout rendre visible pour protéger la
démocratie et lutter contre la corruption, contre le droit fondamental à l'intimité.
Dans l'espace numérique, onde de traçabilité et d'indexation, la transparence
fonctionne comme remède et poison : elle améliore la reddition de comptes
mais fragilise la sphère privée, creuse les asymétries d'information et
expose les individus à des formes nouvelles de violence symbolique.
Considérée sous l'angle de l'investigation numérique, cette tension se traduit
par un conflit entre deux impératifs légitimes : la recherche de preuves et la
protection des données personnelles.

Un cas paradigmatique est l'accès aux logs d'un service public dans le cadre
d'une enquête sur la corruption. L'ouverture de ces traces permet de
reconsitituer les chaînes d'événements, d'identifier des acteurs et de
garantir la transparence administrative. En revanche, la divulgation de ces
logs risque d'exposer des citoyens innocents, de révéler des habitudes
sensible ou des éléments protégés par le secret professionnel. Le risque n'est
pas seulement individuel : l'obsession de la traçabilité érode la confiance,
puisque les usagers modifient leurs comportements pour éviter d'être
surveillés, limitant l'accès à l'information et appauvrissant l'espace public.

L'éthique kantienne offre un cadrage normatif pertinent. Le principe moral
fondamental kantien — agir de telle sorte que l'humanité soit toujours
traitée comme une fin et jamais seulement comme un moyen — impose
une contrainte sur l'usage de la transparence. Appliqué à l'investigation
numérique, ce principe exige que la collecte, la conservation et la divulgation
des preuves ne traitent jamais la personne comme simple ressource
informationnelle. Concrètement, cela se traduit par des règles de proportionnalité,
de minimisation des données, d'anonymisation lorsque possible, et d'un
contrôle judiciaire effectif pour autoriser l'accès aux informations sensibles.

Une résolution pratique combine plusieurs leviers : (i) \emph{principe de
minimisation} : ne collecter que ce qui est strictement nécessaire ; (ii)
\emph{garde-fou procédural} : autorisation judiciaire préalable pour les
accès intrusifs ; (iii) \emph{techniques de protection} : chiffrement, zoning,
masquage et protocoles zero-knowledge pour prouver sans révéler ; (iv)
\emph{transparence procédurale} : audit des accès et droit d'information
pour les personnes affectées. Ainsi la recherche de la vérité reste possible
sans sacrifier la dignité humaine.

En conclusion, la solution n'est ni l'absolue transparence ni le secret total,
mais un équilibre institutionnalisé : une transparence encadrée, respectueuse
des droits, qui fait de l'investigateur non un voyeur mais un gardien de la
déontologie. Ce compromis incarne la lecture kantienne : protéger (la société)
tout en respectant (l'individu) comme une fin en soi.
  \section*{Exercice 2 — Transformation ontologique}
    \subsection*{Comparaison Heidegger / ère numérique}
Heidegger voit la technique comme un mode de révélation du monde (Gestell)
qui transforme la relation humaine au réel. À l'ère numérique, cette idée
se réitère : la technique ne se contente pas d'outiller l'homme ; elle
constitue des modes d'être nouveaux, où l'existence inclut une présence
numérique persistante (profils, historiques, logs). Heidegger insiste sur
l'appauvrissement possible de l'expérience humaine par la réduction du
monde à ressources ; en numérique, cet appauvrissement prend la forme
d'un monde indexé et monétisé.

\subsection*{Analyse d'un profil social comme \emph{être-par-la-trace}}
Un profil social (publications, likes, métadonnées) est une manifestation
intentionalisée de l'individu. Il révèle des choix, habitudes et relations.
Interprété herméneutiquement, il devient un texte : les parties (posts) ne
prennent sens qu'à travers le tout (parcours de vie numérique). L'\emph{être-par-la-trace}
signifie que l'identité s'exprime et se légitime par ces traces, mais qu'elle
n'est pas réductible à elles.

\subsection*{Impact sur la preuve légale}
Les preuves numériques sont immatérielles, volatiles et manipulables.
La transformation ontologique introduit trois défis :
\begin{itemize}
  \item \textbf{Stabilité} : copies binaires et versions multiples ;
  \item \textbf{Authenticité} : garantir origine et intégrité (hash, chain of custody) ;
  \item \textbf{Interprétation} : nécessité d'une herméneutique technique pour éviter
  des conclusions hâtives (contexte, corrélations vs causalité).
\end{itemize} 
  PARTIE 2: 
   \section{Exercice 3} 
   \paragraph{}Nous allons calculer l'entropie de chaque fichier 
\begin{lstlisting}
import math
from collections import Counter
def entropy(filepath):
    with open(filepath, 'rb') as f:
        data = f.read()
    if not data:
        return 0.0
    counts = Counter(data)
    n = len(data)
    H = 0.0
    for c in counts.values():
        p = c / n
        H -= p * math.log2(p)
    return H  # bits par octet

# Usage
 print(entropy('6hearty.txt'))
 print(entropy('OIP (2).jpeg'))
 print(entropy('document.pdf.aes'))
\end{lstlisting}
Voici les résultats obtenus respectivement: 
3.9362600275315254,7.96551329191923,5.9999426266173135
D'après les attentes : on peut conclure que: 
\begin{item}
    \item -Le fichier texte est conforme car son entropie est proche de la borne haute 
    \item -Le fichier image légèrement hors de la borne ne présente rien d'anormal car certains JPEG très compressés comme dans notre cas peuvent atteindre cette valeur 
    \item -le résultat obtenu (très faible) suggère que peut contenir beaucoup d'en-tête non chifffrées par exemple.
\end{item}
Le seuil proposé H<7.7. Si ce n'est pas le cas on peut déduire qu'il s'agit d'un chiffrement suspect ou données chiffrées
\end{document}
